\scanpage{37}
\mspage{14}

\deleted{\unsure{NB}}\quad \uline{Remarques}\quad
1) Avec des hypothèses de ``non sauvagerie'' sur la donnée initiale
\[
 (X,(X_i)_{i\in I},(B_\alpha,\mathcal E_\alpha)_{\alpha\in\Lambda}),
\]
il doit être vrai qu'il existe une structure cônique. Si la stratification
\deleted{\unsure{X}} est équisingulière le long de tout
\[
 X_i^\circ\mathrel{\mathop=^{\mathrm{déf}}}
 X_i\setminus\bigcup_{\substack{j<i\\\unsure{j\le i}}}X_j,
\]
alors il doit \(\unsure{=}\) y avoir une structure cônique équisingulière
(avec la définition ad-hoc évidente) --- il suffit de savoir construire
un tube équisingulier autour des \(X_{I_0}\) \(\cdots\)

2) A priori, si on a une suite \inserted{str.} croissante
de parties
\[
 \varnothing\subset I_0\subset I_1\cdots\subset I_n\cdots\subset I_N=I,
\]
avec \deleted{\(I_i\) formé d'él. minim.}
\inserted{\(\forall i\)} \(I_i\setminus I_{i-1}\) formé d'él. minimaux
de \(I\setminus I_{i-1}\), alors on pourrait définir
\inserted{par récurrence} une notion de structure cônique
pour cette ``filtration'' de \(I\). Mais on
peut montrer que cette notion ne dépend pas
en fait, essentiellement, de la filtration
choisie.

3) Au lieu d'une famille des \((B_\alpha,\mathcal E_\alpha)_{\alpha\in\Lambda}\),
on peut supposer donnée sur \(X\) une
stratification cylindrique \deleted{de type} (le cas envisagé
est celui où l'ens ordonné qui \deleted{\unsure{porte}}
\unsure{décrit} la stratification cylindrique est
\(\mathcal P(\Lambda)^\circ\).

\margin{36}
